%
% einleitung.tex -- Beispiel-File für die Einleitung
%
% (c) 2020 Prof Dr Andreas Müller, Hochschule Rapperswil
%
% !TEX root = ../../buch.tex
% !TEX encoding = UTF-8
%
\section{Die Zeta-Funktion}
\label{basel:section:zeta}

\subsection{Definition der Zeta-Funktion}
\label{basel:subsection:zetadef}

Die \emph{Riemannsche Zeta-Funktion} ist definiert als
\begin{equation}
    \zeta(n)
    =
    \sum_{r=1}^{\infty} \frac{1}{r^n}
    =
    1 + \frac{1}{2^n} + \frac{1}{3^n} + \cdots.
    \label{basel:equation:zetadef}
\end{equation}

\subsection{Konvergenzverhalten der Zeta-Funktion}
\label{basel:subsection:zetakonvergenz}

Wenn $n = 1$ gilt, so divergiert $\zeta(n)$ (sehr langsam).
Diese Reihe ist auch bekannt unter dem Namen \emph{harmonische Reihe}.
Für alle Werte $n > 1$ konvergiert die Zeta-Funktion.

Der Beweis:
Die Terme werden in Zweiergruppen zusammengefasst und mit einer geometrischen
Reihe abgeschätzt:
\begin{align*}
    \sum_{r=1}^{\infty} \frac{1}{r^n}
    &=
    1 + \frac{1}{2^n} + \frac{1}{3^n} + \frac{1}{4^n} + \cdots
    \\
    &=
    1
    + \left(\frac{1}{2^n} + \frac{1}{3^n}\right)
    + \left(\frac{1}{4^n} + \frac{1}{5^n} + \frac{1}{6^n} + \frac{1}{7^n}\right)
    + \cdots
    \\
    &<
    1 + \frac{2}{2^n} + \frac{4}{4^n} + \cdots
    \\
    &=
    1
    + \frac{1}{2^{n-1}}
    + \left(\frac{1}{2^{n-1}}\right)^{\!2}
    + \left(\frac{1}{2^{n-1}}\right)^{\!3}
    + \cdots
    \\
    &=
    \frac{1}{1 - \dfrac{1}{2^{n-1}}}.
\end{align*}
Da die Reihe mit der unendlichen geometrischen Zahlenreihe in einer Ungleichung
in Beziehung gebracht wurde, erhalt man als Grenzwert
\[
    S
    =
    \lim_{n \to \infty} s_n
    =
    \frac{a_1}{1-q},
    \qquad
    \text{mit } q = \frac{1}{2^{n-1}}
    \text{ und } |q| < 1.
\]
Damit die Reihe auch konvergieren kann, muss also gelten:
\[
    \left| \frac{1}{2^{n-1}} \right| < 1
    \;\Rightarrow\;
    |2^{n-1}| > 1
    \;\Rightarrow\;
    n - 1 > 0
    \;\Rightarrow\;
    n > 1.
\]

\subsection{Primzahlen der Zeta-Funktion}
\label{basel:subsection:zetaprimzahlen}

Jede positive ganze Zahl lässt sich als Produkt von Primzahlen ausdrücken.
Daraus folgerte Euler, dass jede ganze positive Zahl $r$ sich für irgendwelche
$r_1, r_2, r_3, \ldots \in \{0, 1, 2, 3, 5, \ldots\}$ in der Form
$r = 2^{r_1} \cdot 3^{r_2} \cdot 5^{r_3}$ ausdrücken lässt, und deswegen gilt
\[
    \frac{1}{r^x}
    =
    \frac{1}{\bigl(2^{r_1} 3^{r_2} 5^{r_3} \cdots\bigr)^x}
    =
    \frac{1}{2^{x r_1} 3^{x r_2} 5^{x r_3} \cdots}.
\]
Fur $x > 1$ ergibt sich daher
\[
    \zeta(x)
    =
    \sum_{r=1}^{\infty} \frac{1}{r^x}
    =
    \sum_{\substack{r_1,\, r_2,\, r_3,\, \ldots\, \geq\, 0}}
        \frac{1}{2^{x r_1} 3^{x r_2} 5^{x r_3} \cdots}
    =
    \left(\sum_{r_1 \geq 0} \frac{1}{2^{x r_1}}\right)
    \left(\sum_{r_2 \geq 0} \frac{1}{3^{x r_2}}\right)
    \left(\sum_{r_3 \geq 0} \frac{1}{5^{x r_3}}\right)
    =
    \prod_{p \text{ prim}}
    \left(\sum_{\alpha=0}^{\infty} \frac{1}{p^{x\alpha}}\right).
\]
Auch hier ist jeder Faktor eine geometrische Reihe und deren Summe beträgt
\[
    \frac{1}{1 - \dfrac{1}{p^x}}
    =
    \frac{1}{1 - p^{-x}}.
\]
Dadurch erhalten wir die \emph{Eulersche Produktformel}:
\begin{equation}
    \zeta(x)
    =
    \sum_{r=1}^{\infty} \frac{1}{r^x}
    =
    \prod_{p \text{ prim}} \frac{1}{1 - p^{-x}},
    \qquad x > 1.
    \label{basel:equation:eulerprod}
\end{equation}
